\chapter{Variáveis Instrumentais}
\label{cap:iv}

\section{O problema de endogeneidade}

OLS produz estimativas viesadas quando o regressor de interesse é correlacionado com
o termo de erro — situação chamada de endogeneidade. As causas mais comuns são:

\begin{itemize}
  \item Variável omitida correlacionada com o regressor
  \item Causalidade reversa (simultaneidade)
  \item Erro de medida no regressor
\end{itemize}

A solução de Variáveis Instrumentais (IV) usa um instrumento $z$ que satisfaz duas
condições: \textbf{relevância} ($z$ correlacionado com o regressor endógeno) e
\textbf{exogeneidade} ($z$ não correlacionado com o erro).

\section{Sintaxe}

\begin{lstlisting}
iv(fórmula_estrutural, fórmula_instrumentos, df)
iv(fórmula_estrutural, fórmula_instrumentos, df, cov=tipo)
\end{lstlisting}

A \textit{fórmula estrutural} especifica a equação de interesse. A
\textit{fórmula dos instrumentos} lista os instrumentos excluídos após \texttt{\textasciitilde}:

\begin{lstlisting}[caption={IV básico — um instrumento}]
// equação estrutural: y ~ x  (x é endógeno)
// instrumento excluído: z
iv(y ~ x, ~ z, df)
\end{lstlisting}

Para múltiplos regressores e instrumentos:

\begin{lstlisting}
iv(y ~ x1 + x2, ~ z1 + z2, df)
\end{lstlisting}

\begin{nota}
  O modelo é identificado quando o número de instrumentos excluídos for maior ou
  igual ao número de regressores endógenos. Com exatamente um instrumento para um
  regressor endógeno, o estimador IV é equivalente ao 2SLS.
\end{nota}

\section{Erros padrão}

As mesmas opções de \lstinline{cov=} disponíveis no OLS aplicam-se ao IV:

\begin{lstlisting}
iv(y ~ x, ~ z, df)                           // clássico
iv(y ~ x, ~ z, df, cov=robust)              // HC1
iv(y ~ x, ~ z, df, cov=HC3)                 // HC3
iv(y ~ x, ~ z, df, cov=cluster, cluster=id) // clusterizado
\end{lstlisting}

\section{Capturando o resultado}

\begin{lstlisting}
let m_iv = iv(y ~ x, ~ z, df)
\end{lstlisting}

O resultado armazenado pode ser usado em \lstinline{esttab} ao lado de modelos OLS:

\begin{lstlisting}[caption={OLS versus IV na mesma tabela}]
let m_ols = ols(y ~ x, df)
let m_iv  = iv(y ~ x, ~ z, df, cov=robust)

esttab(m_ols, m_iv)
\end{lstlisting}

\section{Diagnóstico de instrumentos fracos}

O teste de instrumentos fracos avalia a relevância dos instrumentos na primeira etapa.
Instrumentos fracos inflam o viés IV e invalidam a inferência.

\begin{lstlisting}[caption={Teste de instrumentos fracos}]
weak_iv(y ~ x, ~ z, df)
\end{lstlisting}

A função reporta:
\begin{itemize}
  \item Estatística F da primeira etapa (por variável endógena)
  \item Estatística de Cragg-Donald
  \item Valores críticos de Stock \& Yogo (2005) para $L = 1\ldots10$ instrumentos
  \item Regra de bolso de Staiger \& Stock (1997): $F > 10$
\end{itemize}

\begin{nota}
  A mesma sintaxe de fórmulas de \lstinline{iv()} é usada em \lstinline{weak\_iv()}.
\end{nota}

\section{Valores ajustados}

\begin{lstlisting}
let m_iv = iv(y ~ x, ~ z, df)
let yhat = predict(m_iv, "xb")    // valores ajustados da equação estrutural
\end{lstlisting}

\section{Tabela comparativa}

O fluxo típico de uma análise IV compara OLS e IV para quantificar o viés de
endogeneidade:

\begin{lstlisting}[caption={Fluxo completo de análise IV}]
load "dados.csv" as df

let m_ols = ols(salario ~ educ, df)
let m_iv  = iv(salario ~ educ, ~ distancia, df, cov=robust)

// diagnóstico da primeira etapa
weak_iv(salario ~ educ, ~ distancia, df)

// tabela comparativa
esttab(m_ols, m_iv)
\end{lstlisting}
